% Generated draft: requires mathematical and source review.
\documentclass[11pt]{article}
\usepackage[margin=25mm]{geometry}
\usepackage{fontspec}
\setmainfont{Latin Modern Roman}
\usepackage{amsmath,amssymb,mathtools,unicode-math}
\setmathfont{Latin Modern Math}
\usepackage{enumitem,needspace,xurl,hyperref}
\setlength{\parindent}{0pt}
\setlength{\parskip}{5pt}
\setlength{\emergencystretch}{3em}
\newcommand{\N}{\mathbb{N}}
\newcommand{\Z}{\mathbb{Z}}
\newcommand{\Q}{\mathbb{Q}}
\newcommand{\R}{\mathbb{R}}
\newcommand{\C}{\mathbb{C}}
\newcommand{\F}{\mathbb{F}}
\newcommand{\A}{\mathbb{A}}
\newcommand{\PP}{\mathbb{P}}
\newcommand{\E}{\mathbb{E}}
\newcommand{\eps}{\varepsilon}
\newcommand{\dd}{\,\mathrm{d}}
\begin{document}
\section*{UnsolvedMath ID: 20001648}
\textbf{Source problem code:} catalogue entry
\subsection{Definitions and mathematical statement}
Let \(n\geq2\). The source asks whether there exists ``\(p\to0\)'' such that
\[
 \operatorname{Out}(F_n)/
 \langle\!\langle\gamma^p:\gamma\in\operatorname{Out}(F_n)\rangle\!\rangle
\]
is infinite. The supplied verification reports that \(p\to0\), the exponent and the scope ``all \(\gamma\)'' are present in the PDF, not OCR substitutions. The existential expression \(p\to0\) is not meaningful for an integer exponent. Taking \(p=0\) imposes no relations and does not repair the intended question.

The aid's contextually coherent reconstruction asks for a positive integer exponent \(p\), not necessarily prime, for which the universal exponent quotient
\[
 B_p(G)=G/\langle\!\langle g^p:g\in G\rangle\!\rangle,
 \qquad G=\operatorname{Out}(F_n),
\]
is infinite. Every element of this quotient has order dividing \(p\). It is essential to impose all power relations simultaneously rather than take the quotient by one chosen element's power.

The aid reports a parity obstruction: \(\operatorname{Out}(F_n)\) is generated by involutions, and
\[
 B_p(\operatorname{Out}(F_n))=1\quad\text{for every odd positive integer }p.
\]
Thus an infinite quotient in this reconstruction needs an even exponent; no prime-exponent assertion is supplied.

The positive transfer criterion is the following. Let \(H\trianglelefteq G\) have finite index, suppose \(H\twoheadrightarrow F_r\), \(r\geq2\), and let \(d=\exp(G/H)\), the least common multiple of orders in the finite quotient. If \(B_e(F_r)\) is infinite, then the aid reports
\[
 B_{de}(G)\text{ is infinite}.
\]
In particular sufficiently large odd free-Burnside exponents \(e\) yield infinite quotients of exponent dividing \(de\). The summary applies this to the reported largeness of \(\operatorname{Out}(F_2)\) and \(\operatorname{Out}(F_3)\), where largeness means a finite-index subgroup surjects onto a nonabelian free group. It concludes that infinitely many even exponents work in ranks two and three. It supplies no explicit finite-index subgroup, value of \(d\), or odd-exponent threshold. The reconstructed rank-\(n\geq4\) question is not resolved by this argument.

The original remark adds two separate points. It says DGO implies that \(\operatorname{Out}(F_n)\) is SQ-universal, meaning that every countable group embeds in some quotient; it gives no rank restrictions or complete DGO reference. SQ-universality alone does not specify a universal exponent quotient. It also raises uniform exponential growth for exponentially growing subgroups, and possible uniformity independent of the subgroup. In word-growth terminology these are distinct questions: for each finitely generated subgroup \(L\) of exponential growth, is \(\inf_S\lim_{r\to\infty}|B_S(r)|^{1/r}>1\), and can the lower bound be chosen independently of \(L\) in the ambient group? The intended meaning of exponentially growing subgroup is not specified, so these remain attributed ancillary questions, not conclusions of the Burnside argument.
\subsection{Short English statement}
The source genuinely prints an ill-formed exponent condition. Under the positive-integer reconstruction, the aid rules out every odd exponent and transfers infinite free-Burnside quotients through finite-index free quotients. This gives infinitely many even exponents for \(\operatorname{Out}(F_2)\) and \(\operatorname{Out}(F_3)\), while leaving ranks at least four unresolved in its account.
\subsection{Sources}
\begin{itemize}
\item Local UnsolvedMath registry, permanent ID \(20001648\), code catalogue entry; canonical \texttt{aim-geometric-group-theory-notes.json}, \texttt{notes[156]}.
\item AIM workshop report ``The geometry of the outer automorphism group of a free group,'' Problem 3.6, \url{https://aimath.org/WWN/outerauto/outerauto.pdf}. The supplied verification explicitly reports the PDF's \(p\to0\) and all-power normal closure; this draft does not claim PDF inspection. The verification excerpt ends during its discussion of \(p=0\).
\item The original remark cites ``DGO'' for SQ-universality without an expanded or complete bibliographic reference. The aid also invokes large rank-two/rank-three outer automorphism groups and infinite odd-exponent free Burnside groups, with no complete citations or thresholds supplied.
\item Research attempt 1: result ``partial\_result,'' status ``partially\_solved,'' rigor label ``proved,'' novelty confidence ``low.'' The aid is machine generated, not peer reviewed; independently verify before citation.
\end{itemize}
\end{document}
